% Diapositiva obligatoria: el mejor modelo (igual en los dos decks)
\begin{frame}{El mejor modelo en una diapositiva}
\mensaje{Clasificamos con un índice de ingreso construido desde las personas}
\begin{columns}[T]
\begin{column}{0.6\textwidth}
\resizebox{\linewidth}{!}{% Diagrama del mejor modelo (se usa en los dos decks)
\begin{tikzpicture}[
  font=\scriptsize, >={Stealth[length=2mm]},
  caja/.style={draw=#1, fill=#1!10, rounded corners=3pt, align=center, minimum height=8mm, inner sep=3pt, line width=0.7pt},
  caja/.default=azul, flecha/.style={->, line width=0.7pt, color=gris}]
  % Personas
  \node[caja=naranja, text width=24mm] (per) {\textbf{541 mil personas}\\variables propias\\+ de su hogar};
  \node[caja=naranja, text width=27mm, right=8mm of per, yshift=9mm] (m17) {LightGBM: ingreso total};
  \node[caja=naranja, text width=27mm, right=8mm of per] (m20) {LightGBM: + columnas\\crudas y pesos};
  \node[caja=naranja, text width=27mm, right=8mm of per, yshift=-9mm] (m22) {5 LightGBM: uno por\\componente del ingreso};
  \node[caja=naranja, text width=22mm, right=7mm of m20] (suma) {Suma por hogar\\y promedio: $\hat y$};
  \node[caja=naranja, text width=28mm, right=7mm of suma] (z) {$z = \hat y - g$\\$g$: brecha exacta del DANE\\$\sigma$: error esperado};
  \node[caja=naranja, text width=20mm, below=6mm of z] (logit) {Logit en\\$z,\ z/\sigma,\ \log\sigma$};
  % Hogar
  \node[caja=azul, text width=24mm, below=17mm of per] (hog) {\textbf{165 mil hogares}\\85 variables};
  \node[caja=azul, text width=27mm, right=8mm of hog] (xgb) {XGBoost de\\clasificación};
  % Ensamble
  \node[caja=rosa, text width=34mm, below=6mm of logit] (ens) {$P = 0.3\,P_{\text{XGB}} + 0.7\,P_{\text{ingreso}}$\\pobre si $P \geq 0.40$};
  \draw[flecha] (per.east) -- ++(3mm,0) |- (m17.west);
  \draw[flecha] (per.east) -- (m20.west);
  \draw[flecha] (per.east) -- ++(3mm,0) |- (m22.west);
  \draw[flecha] (m17.east) -- ++(3mm,0) |- (suma.west);
  \draw[flecha] (m20.east) -- (suma.west);
  \draw[flecha] (m22.east) -- ++(3mm,0) |- (suma.west);
  \draw[flecha] (suma) -- (z);
  \draw[flecha] (z) -- (logit);
  \draw[flecha] (logit) -- (ens);
  \draw[flecha] (hog) -- (xgb);
  \draw[flecha] (xgb.east) -- (xgb.east -| ens.west);
\end{tikzpicture}
}
\par\vspace{0.5em}
{\scriptsize\textbf{Especificación (boosting):} $P = 0.3\,P_{\text{XGBoost}} + 0.7\,P_{\text{ingreso}}$. El XGBoost clasifica con 85 variables del hogar. $P_{\text{ingreso}}$: tres LightGBM predicen el ingreso de cada persona; se suma por hogar, se promedian y se compara con la brecha del hogar; un logit lo convierte en probabilidad según qué tan incierta es la predicción ($\sigma$). Pobre si $P \geq 0.40$.\par}
\end{column}
\begin{column}{0.38\textwidth}
\centering
\cifra{0.77}{Kaggle público}\hspace{0.6em}\cifra{0.77}{Kaggle privado}\hspace{0.6em}\cifra{0.776}{F1 CV}
\par\vspace{0.6em}\raggedright
{\footnotesize
\begin{enumerate}\setlength\itemsep{0.25em}
\item \alert{Predecir el ingreso de cada persona} vale más que cambiar de algoritmo.
\item El desbalance se maneja con el \alert{umbral}, no rebalanceando.
\item Sin el módulo de ingresos, una encuesta corta cae a \alert{F1 0.69}.
\end{enumerate}}
\par\vspace{0.3em}
{\tiny Cuatro envíos empatan en 0.77 en el público (Kaggle muestra 2 decimales); entre ellos, el CV difiere en menos de 0.001, menos que su error estándar. Presentamos el que el equipo seleccionó.\par}
\end{column}
\end{columns}
\end{frame}
